\documentclass[10pt]{article}
\usepackage[margin=1in]{geometry}
\usepackage{amsmath,amssymb,amsthm,booktabs}
\usepackage[colorlinks=true,urlcolor=blue,linkcolor=blue]{hyperref}
\newtheorem{theorem}{Theorem}
\title{A counterexample to the higher Toeplitz\\
Hurwitz-stability conjecture\\\large MathDB 323159}
\author{Principia Math In Progress}
\date{11 September 2026}
\begin{document}
\maketitle

\section*{Task 1. Resolve the stability conjecture}

Karp's Conjecture 5~\cite{K} asserts that the rising-factorial determinant
polynomial $P_n^r(x)$ is Hurwitz stable whenever the input sequence belongs
to $PF_r$, for $r\geq2$. Hurwitz stability means that every zero has
strictly negative real part.
The order-two specialization, explicitly included in the source, is
\begin{equation}
 P_n^2(x)=\sum_{k=0}^n\binom nk f_kf_{n-k}
 \bigl[(x)_k(x)_{n-k}-(x+1)_k(x-1)_{n-k}\bigr],
 \label{eq:definition}
\end{equation}
where $(x)_k$ is the rising factorial.

\begin{theorem}
Take $r=2$, $n=15$ and
\[
 (f_0,\ldots,f_{15})=
 (1,51,52,53,54,55,56,57,57,56,55,54,53,52,51,1).
\]
This is a positive strictly log-concave $PF_2$ sequence, but $P_{15}^2$
has exactly one zero in
\[
 \left|x-\frac{1+41i}{10}\right|<\frac1{20}.
\]
In particular it has a zero with $\operatorname{Re}x>1/20$,
and Conjecture 5 is false.
\end{theorem}

\section*{The input satisfies the full hypothesis}

The adjacent log-concavity differences are
\[
 f_k^2-f_{k-1}f_{k+1}=
 \begin{cases}
 2549,&k=1,14,\\
 57,&k=7,8,\\
 1,&k=2,\ldots,6,\ 9,\ldots,13.
 \end{cases}
\]
Thus strict log-concavity holds at every required index.

For completeness, extend $f_k$ by zero outside $0\leq k\leq15$.
All entries of its infinite Toeplitz matrix are nonnegative.
A two-by-two minor has the form
\[
 f_a f_{a+h-d}-f_{a+h}f_{a-d},\qquad h,d>0
\]
for integer indices. If the second product is zero, the minor is
nonnegative. Otherwise all indices from $a-d$ through $a+h$
are in the positive support. Since the ratios $f_j/f_{j-1}$ decrease,
\[
 \frac{f_a}{f_{a-d}}
 \geq\frac{f_{a+h}}{f_{a+h-d}},
\]
by multiplying $d$ successive ratio comparisons.
The minor is again nonnegative. This proves $PF_2$ directly, including
every minor of the infinite matrix.

\section*{An explicit polynomial factorization}

Expanding the finite sum~\eqref{eq:definition} gives
\begin{equation}
 P_{15}^2(x)=210(x+1)_7 R(x),
 \label{eq:factor}
\end{equation}
where
\begin{align*}
 R(x)={}&191647x^6+4104969x^5+36750955x^4+192645735x^3\\
        &+736979998x^2+2128319736x+3212352000.
\end{align*}
This is an integer polynomial identity. It can be checked by ordinary
coefficient multiplication. The accompanying verifier also checks it by
evaluating both sides independently at $x=0,1,\ldots,15$;
both sides have degree at most $15$, so these sixteen exact evaluations
alone certify the identity.

All coefficients of the displayed factorization are positive.
Thus the counterexample concerns stability; it does not contradict
coefficient positivity.

\section*{An integer certificate for a right-half-plane zero}

Put
\[
 T(w)=20^6 R\left(\frac{2+82i+w}{20}\right)
      =\sum_{j=0}^6 c_jw^j.
\]
Its coefficients are the following Gaussian integers:
\begin{center}
\small
\begin{tabular}{r rr}
\toprule
$j$&$\operatorname{Re}c_j$&$\operatorname{Im}c_j$\\
\midrule
0&6242984901271040&-4815010659523584\\
1&-7146943743914496&-9080441566052096\\
2&-369163407187520&-56477312908800\\
3&-4012914867840&2981422413440\\
4&-3796641800&34603649040\\
5&84399144&94290324\\
6&191647&0\\
\bottomrule
\end{tabular}
\end{center}
Let
\[
 L=11555657658184807,\qquad
 U=\sum_{j\ne1}\left\lceil
       \sqrt{(\operatorname{Re}c_j)^2+(\operatorname{Im}c_j)^2}
       \right\rceil
   =8262602600766508 .
\]
Direct integer squaring verifies $L^2\leq |c_1|^2$, and integer square
root bounds give the displayed exact value of $U$.
Moreover,
\[
 L-U=3293055057418299>0.
\]
Consequently on $|w|=1$,
\[
 |T(w)-c_1w|
 \leq\sum_{j\ne1}|c_j|
 \leq U<L\leq|c_1w|.
\]
Rouch\'e's theorem implies that $T$ and $c_1w$ have the same number
of zeros inside the unit disk, namely one, counted with multiplicity.
Undoing the affine change of variable gives exactly one zero of $R$
in the claimed disk. The factor $(x+1)_7$ has its zeros at
$-1,-2,\ldots,-7$, outside this disk, so the same statement holds for
$P_{15}^2$. Every point in the disk has real part greater than $1/20$.
This completes the refutation.

\section*{Verification and scope}

The verifier uses only exact integer arithmetic and Python's standard
library. It checks the input inequalities, constructs the original
determinant polynomial, checks~\eqref{eq:factor} by coefficients and by
independent evaluations, and computes $T$ by both binomial expansion
and Horner substitution. The norm comparisons use integer square roots.
No approximate root finder is used in the certificate.

The result refutes the universal assertion through its allowed $r=2$
case. It does not decide the separate assertions for each $r\geq3$,
or the stronger $PF_\infty$ real-root conjecture.
The source's previously known nonreal-zero examples under mere
log-concavity do not assert failure of Hurwitz stability.
The published statement, live MathDB status and public claims were checked
on 11 September 2026; no matching existing stability-refutation claim
was located. Details are in \texttt{literature-audit.md}.

\section*{Tasks 2--5}

\begin{enumerate}
\setcounter{enumi}{1}
\item Determine the smallest degree index $n$ for which positive strictly
log-concave data give an unstable order-two determinant polynomial.
\item Determine the instability ranges in $A$ and $n$ for
$f_0=f_n=1$ and $f_k=A+\min(k,n-k)$ at the interior indices.
\item Prove or refute Hurwitz stability separately for each determinant
order $r\geq3$ under the corresponding $PF_r$ hypothesis.
\item Resolve the proposed real-negative-zero property under the
stronger $PF_\infty$ hypothesis.
\end{enumerate}

\begin{thebibliography}{9}
\bibitem{K}
D.~Karp, \emph{Positivity of Toeplitz determinants formed by rising
factorial series and properties of related polynomials},
Zap. Nauchn. Sem. POMI \textbf{404} (2012), 184--198,
Conjecture 5 and the order-two identity, p.~193.
\href{https://www.mathnet.ru/eng/znsl5267}{Published article};
\href{https://arxiv.org/abs/1203.1482}{arXiv:1203.1482}.
English translation: Journal of Mathematical Sciences
\textbf{193} (2013), 106--114,
\href{https://doi.org/10.1007/s10958-013-1438-y}
{doi:10.1007/s10958-013-1438-y}.
\end{thebibliography}
\end{document}
